
\section{Transparency}\fbeg{}\begin{frame}{\ft{Data Transparency Becomes a Priority}}

\vspace{-.5em}	

{\Large\fontfamily{uhv}\selectfont

\hspace*{20pt}\begin{minipage}{.96\textwidth}
\vspace{4pt}
		
\begin{lightquadblockc}{0,0.1,0.1}{\parbox{21cm}{\centering \vspace{3pt}}}
\begin{center}\begin{minipage}{\textwidth}
{\Huge \setlength{\leftmargini}{30pt}\begin{enumerate}
\dmitem {\dsp} Publishers may require authors to share raw data or else explain specifically 
why it is not possible to do so 
\vspace{20pt}
\dmitem {\dsp} Funding agencies may require data transparency as a precondition for 
grantees (e.g., the Bill and Melinda Gates Foundation)
\vspace{20pt}
\dmitem {\dsp} \q{Data Availability} or \q{Supplemental Materials} sections are now common on 
publication landing pages
\vspace{20pt}
%\footnotemark[2]\footnotetext{\Large \raisebox{7pt}{\em{2}} 
%	\hspace{7pt}See Slides 11 and 12 for details.}
\end{enumerate}
}\end{minipage}
\end{center}
\end{lightquadblockc}
\end{minipage}

}

\vspace{-11pt}
{\thrule}
\hspace{32pt}\begin{minipage}{.8\textwidth}
{\fontsize{19}{25}\selectfont \setlength{\leftmargini}{3pt}
Authors however often confuse tables or summaries with actual data sharing --- %\vspace{.2em} 
i.e., publishing raw data sets (not just a statistical overview)

}
\end{minipage}


\end{frame}

